\documentclass[11pt,a4paper]{article}

\usepackage[utf8]{inputenc}
\usepackage[T1]{fontenc}
\usepackage[margin=2.5cm]{geometry}
\usepackage{amsmath,amssymb,amsthm}
\usepackage{booktabs}
\usepackage{microtype}
\usepackage{xcolor}
\usepackage{listings}
\usepackage[backend=biber,style=numeric,sorting=nyt,maxbibnames=99]{biblatex}
\usepackage[hidelinks]{hyperref}

\addbibresource{efx.bib}

% ---- theorem environments --------------------------------------------------
\theoremstyle{definition}
\newtheorem{definition}{Definition}
\theoremstyle{plain}
\newtheorem{theorem}{Theorem}
\newtheorem{corollary}{Corollary}
\newtheorem{proposition}{Proposition}
\newtheorem{observation}{Observation}

% ---- listing style ---------------------------------------------------------
\definecolor{codekw}{rgb}{0.10,0.10,0.65}
\definecolor{codecm}{rgb}{0.30,0.45,0.30}
\definecolor{codebg}{rgb}{0.97,0.97,0.97}
\lstdefinestyle{code}{
  basicstyle=\ttfamily\footnotesize,
  keywordstyle=\color{codekw}\bfseries,
  commentstyle=\color{codecm}\itshape,
  backgroundcolor=\color{codebg},
  numberstyle=\tiny\color{gray},
  numbers=left,
  numbersep=7pt,
  frame=single,
  framerule=0.3pt,
  rulecolor=\color{gray!50},
  breaklines=true,
  showstringspaces=false,
  columns=fullflexible,
  keepspaces=true,
  tabsize=2,
}
\lstset{style=code}

% ---- macros ----------------------------------------------------------------
\newcommand{\EFX}{\textsf{EFX}}
\newcommand{\EFone}{\textsf{EF1}}
\newcommand{\agents}{N}
\newcommand{\goods}{M}
\newcommand{\val}[1]{v_{#1}}

\title{\textbf{Full \EFX{} Allocations Exist for Four Agents and Eight Goods\\
with Additive Valuations}}
\author{Eduard Tykhoniuk}
\date{June 2026}

\begin{document}
\maketitle

\begin{abstract}
The existence of \emph{envy-free up to any good} (\EFX) allocations is one of the
central open problems in the fair division of indivisible goods. It is known to
hold for two agents with general valuations and for three agents with additive
valuations; for four agents it is known to hold
\emph{almost} fully --- an \EFX{} allocation always exists if one item may be
left unallocated. Whether a \emph{full} (complete) \EFX{} allocation always
exists for four agents remains open. In this report we settle the question
\emph{for a fixed instance size}: we show that every \textbf{additive} profile
with four agents and eight goods admits a full \EFX{} allocation. The argument
reduces the existence statement to a single satisfiability query over linear real
arithmetic: the negation of \EFX{} existence is satisfiable if and only if a
counterexample profile exists. We encode this query with the Z3 SMT solver and
exploit a goods-relabelling symmetry that, as we verify empirically, is needed for
the larger instances to finish; after roughly $62.5$ hours the solver returns
\texttt{unsat}.
To the best of our knowledge this is the first verification of full \EFX{}
existence for four agents and eight goods under additive valuations.
\end{abstract}

\tableofcontents
\medskip
\hrule

\section{Introduction}
\label{sec:intro}

\emph{Fair division} studies how to partition a set of resources among agents
with possibly conflicting preferences so that the outcome is, in some precise
sense, fair. When the resources are \emph{indivisible goods} --- items that
cannot be split, such as houses, courses, tasks, or shifts --- the classical
benchmark of \emph{envy-freeness} (no agent prefers another agent's bundle to
its own) is generally unachievable: with one good and two agents, whoever leaves
empty-handed envies the other. Research has therefore concentrated on
\emph{relaxations} of envy-freeness that are achievable yet still meaningful.

The most prominent relaxations are \emph{envy-free up to one good} (\EFone) and
\emph{envy-free up to any good} (\EFX). Under \EFone{} an agent may envy another,
provided the envy disappears once \emph{some} single good is removed from the
envied bundle; \EFone{} allocations always exist and can be computed in
polynomial time, even for general monotone valuations
\cite{LiptonEtAl2004}; see also \cite{Budish2011}. \EFX{} is the strictly stronger requirement
that the envy disappears once \emph{any} single good is removed
\cite{Caragiannis2019,PlautRoughgarden2020}. \EFX{} is one of the strongest such
fairness notions for indivisible goods, and whether an \EFX{} allocation always
exists is a major open problem in the area \cite{AmanatidisSurvey2023}.

\paragraph{State of the art.}
\EFX{} existence is known in only a handful of regimes. It holds for two agents
with arbitrary monotone valuations \cite{PlautRoughgarden2020}, and for three
agents with additive valuations \cite{ChaudhuryGargMehlhorn2024}. For four
agents, Berger, Cohen, Feldman and Fiat \cite{BergerCohenFeldmanFiat2022} proved
that an \EFX{} allocation always exists \emph{if at most one good is left
unallocated} --- an ``almost full'' guarantee --- and, more generally, that with
$\agents$ additive agents an \EFX{} allocation exists leaving at most
$\agents-2$ goods unallocated. Whether a \emph{full} \EFX{} allocation (one that
allocates \emph{all} goods) always exists for four agents is, to date, open even
for additive valuations. On the negative side, Akrami, Mayorov, Mehlhorn,
Srinivas and Weidenbach \cite{AkramiEtAl2026} very recently produced the first
counterexamples to \EFX{} existence under \emph{submodular} (and hence monotone)
valuations, showing that for $n\geq 3$ agents \EFX{} need not exist once
$m \geq n+5$; they leave the \emph{additive} case wide open.

Beyond exact existence, a large body of work studies what can be guaranteed when
exactness is relaxed: approximate \EFX{} (where an agent's value for its own
bundle is within a multiplicative factor of the envied bundle minus a good),
\EFX{} with bounded charity (a small number of unallocated goods, as above),
restricted valuation classes (binary, two-valued, or few distinct types), and
\EFX{} on graphs where agents only compare with neighbours. The survey of
Amanatidis et al.\ \cite{AmanatidisSurvey2023} gives a comprehensive map of these
directions and frames exact \EFX{} existence for general additive valuations as
a central open question.

\paragraph{This report.}
We approach full \EFX{} existence \emph{computationally}, for one fixed instance
size. Concretely, we prove:

\begin{theorem}[Main result]
\label{thm:main}
Every instance with four agents, eight goods and non-negative additive
valuations admits a full \EFX{} allocation.
\end{theorem}

The proof is computational. The statement ``\EFX{} always exists for a fixed
$(\agents,\goods)$'' has the form: \emph{for every} valuation profile \emph{there
exists} an allocation that is \EFX. Its negation reads: \emph{there exists} a
profile such that \emph{every} allocation fails \EFX. Because the valuations are
real-valued and the \EFX{} conditions are linear inequalities, this negation is a
single existential query in linear real arithmetic, which an SMT solver decides
exactly: a satisfying assignment is a counterexample profile, while
unsatisfiability means no counterexample exists, i.e.\ \EFX{} always exists.
Running the query for $\agents=4,\ \goods=8$ returns \texttt{unsat}, which
establishes Theorem~\ref{thm:main} for strictly positive valuations; a short
reduction (Section~\ref{sec:nonneg}) then removes the positivity assumption,
extending it to all non-negative profiles in the strong \EFX$_0$ sense.

The remainder of the report sets up this reduction
(Section~\ref{sec:reduction}), describes the Z3-based implementation and a
goods-relabelling symmetry that we find empirically to be needed for the larger
instances to finish (Section~\ref{sec:impl}), reports the result
(Section~\ref{sec:results}), and places it against the literature
(Section~\ref{sec:related}).

\section{Preliminaries}
\label{sec:prelim}

\paragraph{Instances.}
There are $\agents$ \emph{agents} $[\agents]=\{1,\dots,\agents\}$ and $\goods$
indivisible \emph{goods} $[\goods]=\{1,\dots,\goods\}$. Each agent $i$ has an
\emph{additive} valuation $\val{i}\colon 2^{[\goods]}\to\mathbb{R}_{\geq 0}$,
meaning $\val{i}(S)=\sum_{g\in S}\val{i}(\{g\})$ for every bundle
$S\subseteq[\goods]$. We write $w_{i,g}:=\val{i}(\{g\})$ for the value agent $i$
assigns to good $g$, so a profile is an $\agents\times\goods$ matrix
$(w_{i,g})$. Throughout we assume valuations are \emph{strictly positive},
$w_{i,g}>0$; this is the standard ``goods'' setting and the one our solver
encodes.

\paragraph{Allocations.}
An \emph{allocation} is an ordered partition $X=(X_1,\dots,X_\agents)$ of a
subset of $[\goods]$ into pairwise-disjoint bundles, where $X_i$ is the bundle
given to agent $i$. The allocation is \emph{full} (or \emph{complete}) if
$\bigcup_i X_i = [\goods]$, i.e.\ no good is left over.

\begin{definition}[Envy notions]
\label{def:efx}
Fix valuations $\val{1},\dots,\val{\agents}$ and an allocation $X$. In each clause
below, $i$ and $j$ range over all ordered pairs of agents.
\begin{itemize}
  \item $X$ is \emph{envy-free} (\textsf{EF}) if
        \[
          \val{i}(X_i)\ \geq\ \val{i}(X_j).
        \]
  \item $X$ is \emph{envy-free up to one good} (\EFone) if
        \[
          \exists\, g\in X_j:\quad
          \val{i}(X_i)\ \geq\ \val{i}(X_j\setminus\{g\})
        \]
        (and is defined to hold when $X_j=\emptyset$).
  \item $X$ is \emph{envy-free up to any good} (\EFX) if
        \[
          \forall\, g\in X_j:\quad
          \val{i}(X_i)\ \geq\ \val{i}\bigl(X_j\setminus\{g\}\bigr).
        \]
\end{itemize}
The only difference between \EFone{} and \EFX{} is the quantifier on $g$ ($\exists$
versus $\forall$): \EFone{} asks that \emph{some} good can be removed to eliminate
the envy, \EFX{} that \emph{every} good does.
\end{definition}

Equivalently, $X$ is \EFX{} iff for every ordered pair $(i,j)$,
\[
  \val{i}(X_i)\ \geq\ \val{i}(X_j) - \min_{g\in X_j}\val{i}(\{g\}),
\]
because under additive (positive) valuations the bundle $X_j\setminus\{g\}$ of
largest value to $i$ is the one obtained by removing $i$'s \emph{least}-valued
good in $X_j$. When $X_j=\emptyset$ there is no good to remove and the condition
holds trivially.

\paragraph{A note on the literature.}
Two variants of \EFX{} appear in the literature \cite{AmanatidisSurvey2023},
differing in how they treat a good the envious agent values at zero. \EFX$_+$, as
originally defined by Caragiannis et al.\ \cite{Caragiannis2019}, removes only
\emph{positively}-valued goods ($g\in X_j$ with $\val{i}(\{g\})>0$); the stronger
\EFX$_0$, used by Plaut and Roughgarden \cite{PlautRoughgarden2020}, removes
\emph{any} good $g\in X_j$, including one valued at zero. The corresponding
existence guarantees are equivalent at any fixed instance size
(Corollary~\ref{cor:equiv}), so the choice between them does not affect our
results. We adopt the stronger \EFX$_0$ notion and state all of our results for
it; the encoding of Section~\ref{sec:reduction} tests precisely this condition.

The existence question we study is:

\begin{quote}
\emph{For fixed $\agents$ and $\goods$, does every additive profile admit a full
\EFX{} allocation?}
\end{quote}

A \emph{counterexample} is a profile for which \emph{no} full \EFX{} allocation
exists. Theorem~\ref{thm:main} asserts that for $(\agents,\goods)=(4,8)$ there is
no counterexample.

\section{Reducing existence to satisfiability}
\label{sec:reduction}

\subsection{Negating the existence statement}

Fix $(\agents,\goods)$ and let $\mathcal{P}$ denote the (finite) set of ordered
partitions of $[\goods]$ into $\agents$ non-empty bundles. The statement we want
to prove is
\[
  \underbrace{\forall\, (w_{i,g})>0}_{\text{every profile}}\quad
  \underbrace{\exists\, P\in\mathcal{P}}_{\text{some allocation}}\quad
  \mathrm{EFX}(P,w),
\]
where $\mathrm{EFX}(P,w)$ is the predicate ``$P$ is \EFX{} under profile $w$''.
There are infinitely many profiles, so the statement cannot be decided by checking
them one by one. Its negation, however, is purely existential in
$w$,
\[
  \exists\, (w_{i,g})>0\quad
  \bigwedge_{P\in\mathcal{P}}\ \neg\,\mathrm{EFX}(P,w).
\tag{$\star$}
\label{eq:negation}
\]
That is: a counterexample is a single profile for which every allocation fails
\EFX. Statement~\eqref{eq:negation} is a conjunction of disjunctions of
\emph{linear} inequalities in the real variables $w_{i,g}$ --- the fragment
\emph{linear real arithmetic} (LRA) that SMT solvers decide
exactly.\footnote{Linearity also bounds the search a priori: if
\eqref{eq:negation} is satisfiable then it has a \emph{rational} solution, because
a feasible system of linear inequalities with rational coefficients has a rational
point (the feasible region of the strict system is open and non-empty, hence
contains a rational point). A counterexample profile, if one exists, may therefore
be taken rational, matching the exact rational arithmetic the solver uses
(Section~\ref{sec:lra}).} Feeding \eqref{eq:negation} to such a solver yields:
\begin{itemize}
  \item \texttt{sat} $\Rightarrow$ a counterexample profile, which can be read off
        as a rational matrix; or
  \item \texttt{unsat} $\Rightarrow$ \eqref{eq:negation} has no solution, hence the
        original statement holds: \EFX{} always exists.
\end{itemize}

\subsection{Deciding the query: linear real arithmetic}
\label{sec:lra}

Why can a solver decide \eqref{eq:negation} \emph{exactly}, rather than only
search heuristically? The formula \eqref{eq:negation} lives in \emph{quantifier-free
linear real arithmetic} (QF-LRA): a Boolean combination of linear inequalities
$\sum_g a_g\,w_{i,g}\bowtie c$ with rational coefficients, where
$\bowtie\,\in\{<,\leq,=,\geq,>\}$. This theory is \emph{decidable}: a
\textbf{conjunction} of linear inequalities (an intersection of half-spaces) over
the reals is satisfiable iff the corresponding polyhedron is non-empty, which a
simplex-style procedure settles in finitely many exact-rational pivots; an SMT
solver lifts this to arbitrary Boolean structure by searching over truth
assignments to the inequalities in a SAT-solving loop (the DPLL($T$) procedure
underlying SMT). All arithmetic is carried out over the rationals $\mathbb{Q}$,
with no floating-point approximation and no $\varepsilon$-tolerances, so an
\texttt{unsat} verdict is a true impossibility rather than a numerical
artefact --- provided the solver reasons exactly. Z3's QF-LRA core is designed to do so, using
exact arbitrary-precision rational arithmetic. This exactness is what lets us treat valuations as
arbitrary positive reals: proving \texttt{unsat} over $\mathbb{R}$ rules out
\emph{all} real --- in particular all rational and all integer --- profiles at
once.

\subsection{Encoding a single allocation: the \texttt{nonenvy} predicate}

For a fixed partition $P=(P_1,\dots,P_\agents)$, the predicate
$\mathrm{EFX}(P,w)$ expands into one inequality per (envious agent, envied agent,
removed good) triple. For agent $i$ looking at agent $j$'s bundle $P_j$, and for
each good $k\in P_j$, \EFX{} demands
\begin{equation}
  \sum_{g\in P_i} w_{i,g}\ \geq\ \sum_{g\in P_j} w_{i,g} - w_{i,k}.
  \label{eq:efxconstraint}
\end{equation}
$\mathrm{EFX}(P,w)$ is exactly the conjunction of \eqref{eq:efxconstraint} over all
$i\neq j$ and all $k\in P_j$. Two simplifications are sound and are used in the
implementation:
\begin{itemize}
  \item \textbf{An envied bundle that is a singleton imposes no constraint}: if
        $|P_j|\leq 1$, the terms from $j$ may be dropped
        (Observation~\ref{obs:singleton}).
  \item \textbf{Own bundle is excluded} ($i=j$ is skipped), since an agent never
        envies itself.
\end{itemize}

\begin{observation}[Singletons impose no \EFX{} constraint]
\label{obs:singleton}
Under non-negative valuations, if $|P_j|\leq 1$ then no agent $i\neq j$ can
\EFX-envy $j$, so the corresponding constraints \eqref{eq:efxconstraint} always
hold and can be dropped.
\end{observation}

\begin{proof}
If $P_j=\emptyset$ there is no good $k$ to quantify over and the predicate is
empty. If $P_j=\{k\}$, the only instance of \eqref{eq:efxconstraint} reads
$\sum_{g\in P_i}w_{i,g}\geq w_{i,k}-w_{i,k}=0$, which holds because all
valuations are non-negative. In both cases the singleton adds no real constraint.
\end{proof}

\paragraph{Example expansion.}
In general a partition contributes one inequality for every ordered pair $(i,j)$
with $|P_j|\geq 2$ and every good $k\in P_j$. For $\agents=2$, $\goods=3$ and the
partition $P=(\{1\},\{2,3\})$ only the pair $(1,2)$ qualifies --- agent~$2$ faces
the singleton $P_1=\{1\}$, which contributes nothing by
Observation~\ref{obs:singleton} --- so the block is small:
\[
  \underbrace{w_{1,1}\ \geq\ (w_{1,2}+w_{1,3})-w_{1,2}}_{\text{remove good }2}
  \quad\wedge\quad
  \underbrace{w_{1,1}\ \geq\ (w_{1,2}+w_{1,3})-w_{1,3}}_{\text{remove good }3},
\]
i.e.\ $w_{1,1}\geq w_{1,3}$ and $w_{1,1}\geq w_{1,2}$. Each partition yields such a
block, and the solver asserts the negation of their disjunction: that no
partition's block holds in full.

\subsection{Why restricting to non-empty partitions is conservative}
\label{sec:conservative}

The set $\mathcal{P}$ ranges only over partitions in which \emph{every} agent
receives at least one good. \EFX{} allocations are not required to be of this
form, so $\mathcal{P}$ omits some valid candidate allocations (those that hand
an agent the empty bundle). It is important to see that this restriction is
\emph{safe for the existence direction} we care about.

\begin{proposition}[Soundness of the restricted search]
\label{prop:conservative}
Let $\mathcal{P}'\subseteq\mathcal{P}^{\mathrm{all}}$ be any subset of the full
set of allocations. If the query
$\exists w>0:\ \bigwedge_{P\in\mathcal{P}'}\neg\mathrm{EFX}(P,w)$ is
unsatisfiable (has no solution $w$, i.e.\ no counterexample), then \EFX{} always
exists (even when ranging over $\mathcal{P}^{\mathrm{all}}$).
\end{proposition}

\begin{proof}
Removing allocations from the conjunction in \eqref{eq:negation} can only
\emph{weaken} it: $\bigwedge_{P\in\mathcal{P}^{\mathrm{all}}}\neg\mathrm{EFX}
\implies \bigwedge_{P\in\mathcal{P}'}\neg\mathrm{EFX}$. Contrapositively, if the
smaller conjunction over $\mathcal{P}'$ is unsatisfiable, so is the larger one
over $\mathcal{P}^{\mathrm{all}}$; hence no counterexample profile exists and
\EFX{} always exists.
\end{proof}

Restricting the candidate set can therefore only introduce false counterexamples,
never false claims of existence; an \texttt{unsat} on the restricted set
$\mathcal{P}$ still proves existence. This is what our run establishes for
$(\agents,\goods)=(4,8)$.\footnote{In any case, for positive additive valuations
with $\goods\geq\agents$ an allocation that leaves an agent empty is never \EFX{}
unless every other bundle is a singleton, so the restriction discards little here;
Proposition~\ref{prop:conservative} makes the argument independent of this.}

\subsection{Symmetry reduction}
\label{sec:symmetry}

Relabelling the goods is a symmetry of the problem: permuting the columns of a
profile $w$ and applying the same permutation to every bundle preserves \EFX. We
break this symmetry by fixing one agent's view. Without loss of generality the
goods are indexed so that agent $1$'s valuations are non-decreasing,
\[
  w_{1,1}\ \leq\ w_{1,2}\ \leq\ \cdots\ \leq\ w_{1,\goods}.
\]
This is encoded as $\goods-1$ extra inequalities. It is sound because every
profile is equivalent, after relabelling goods, to one in sorted form, and it
removes a factor of up to $\goods!$ from the search space (for $\goods=8$, up to
$40320$). The realised speedup can be smaller than $\goods!$, since Z3 already
performs some symmetry reduction internally; Section~\ref{sec:bench} measures the
actual effect.

\subsection{Counting the candidate allocations}

The number of ordered partitions of $\goods$ goods into $\agents$ non-empty
labelled bundles is the surjection count
$\agents!\,\genfrac{\{}{\}}{0pt}{}{\goods}{\agents}$, where
$\genfrac{\{}{\}}{0pt}{}{\goods}{\agents}$ is a Stirling number of the second
kind. For the three sizes we ran ($\agents=4$, $\goods\in\{6,7,8\}$):
\[
  4!\cdot\genfrac{\{}{\}}{0pt}{}{6}{4}=1560,\qquad
  4!\cdot\genfrac{\{}{\}}{0pt}{}{7}{4}=8400,\qquad
  4!\cdot\genfrac{\{}{\}}{0pt}{}{8}{4}=40824.
\]
Each of these $40824$ partitions contributes its block of inequalities
\eqref{eq:efxconstraint}; the solver must refute the conjunction of all their
negations simultaneously.

\subsection{The shape of the search space}
\label{sec:searchspace}

We now explain why the problem is hard, which accounts for both the runtime and
the limits of the approach. The formula \eqref{eq:negation}
has only $32$ real variables, but a large Boolean structure: a counterexample
profile must violate at least one \EFX{} inequality of each of the $40824$
allocations at once. The solver does not go through profiles or
combinations one by one. Instead it uses the linear structure to prune: once a
partial set of constraints has no solution, every larger set that contains it has
no solution either, so the solver can skip them all.

Two consequences follow. First, the difficulty is highly sensitive to $\goods$:
going from $\goods=6$ to $\goods=8$ multiplies the allocation count by about $26$
and lengthens every constraint, and the running time grows from under a second to
about $62$ hours (Section~\ref{sec:results}). Second, this sensitivity also means
the $(4,8)$ run is computationally demanding: the related three-agent /
seven-good \EFX{} query did not terminate within a $40$-hour budget for the same
Z3-over-LRA approach \cite{AkramiEtAl2026}, so termination at this scale cannot be
taken for granted.

\section{Implementation}
\label{sec:impl}

\subsection{Z3 pipeline}

The main prover is a C++ program built against the Z3 C++ API
\cite{deMouraBjorner2008}. It (i) creates the $\agents\times\goods$ real
variables $w_{i,g}$ and asserts positivity and the sortedness constraints of
Section~\ref{sec:symmetry}; (ii) enumerates the non-empty ordered partitions;
(iii) builds, for each partition, the conjunction \texttt{nonenvy} of
constraints \eqref{eq:efxconstraint}; (iv) asserts the negation of their
disjunction; and (v) calls the solver. The core of the encoding is shown below.

\begin{lstlisting}[language=C++,caption={The \texttt{nonenvy} predicate
(constraints \eqref{eq:efxconstraint}), from \texttt{main.cpp}.},label={lst:nonenvy}]
z3::expr nonenvy(z3::context &ctx,
                 const std::vector<std::vector<int>> &partition,
                 const std::vector<std::vector<z3::expr>> &w) {
  z3::expr_vector constraints(ctx);
  for (int i = 0; i < N; ++i) {
    z3::expr total = ctx.real_val(0);          // value of i's own bundle
    for (int j : partition[i]) total = total + w[i][j];

    for (int j = 0; j < N; ++j) {
      if (i == j) continue;
      if (partition[j].size() <= 1) continue;  // singletons: EFX is automatic
      z3::expr other = ctx.real_val(0);        // value of j's bundle to i
      for (int k : partition[j]) other = other + w[i][k];
      // For every k in j's bundle: i must not envy j after removing k.
      for (int k : partition[j])
        constraints.push_back(total >= other - w[i][k]);
    }
  }
  return z3::mk_and(constraints);
}
\end{lstlisting}

\begin{lstlisting}[language=C++,caption={Assembling and refuting the global
disjunction, from \texttt{main.cpp}.},label={lst:assemble}]
// First agent's valuations sorted (goods-relabelling symmetry break).
for (int j = 0; j < M - 1; ++j) solver.add(w[0][j] <= w[0][j + 1]);

auto partitions = ordered_groups(N, M);           // 40824 for (4,8)
z3::expr_vector valid_exprs(ctx);
for (const auto &part : partitions)
  valid_exprs.push_back(nonenvy(ctx, part, w));
z3::expr or_valid = z3::mk_or(valid_exprs);
solver.add(!or_valid);                             // assert (*) : no allocation is EFX

z3::check_result result = solver.check();          // unsat  =>  EFX always exists
\end{lstlisting}

The program also exports the assembled query in the standard SMT-LIB~2 format so
that it can be replayed by any compliant solver; for $(\agents,\goods)=(4,8)$
this script (\texttt{problem\_4\_8.smt2}) is about $9$\,MiB, compared with
$0.24$\,MiB for $(4,6)$. Its structure is a flat declaration of the $32$ real
variables, the positivity and sortedness assertions, and one large negated
disjunction:

\begin{lstlisting}[caption={Shape of the exported SMT-LIB~2 query (abridged).},label={lst:smtlib}]
(declare-fun w_0_0 () Real)  ... (declare-fun w_3_7 () Real)
(assert (> w_0_0 0)) ...                      ; positivity, 32 of these
(assert (<= w_0_0 w_0_1)) ...                 ; agent 0 sorted, 7 of these
(assert (not (or
   (and (>= ...) (>= ...) ...)                ; nonenvy(partition 1)
   ...
   (and (>= ...) (>= ...) ...))))             ; nonenvy(partition 40824)
(check-sat)
\end{lstlisting}

\subsection{The relabelling symmetry break, measured}
\label{sec:bench}

Section~\ref{sec:symmetry} argued that fixing agent~$1$'s valuations in
non-decreasing order removes a factor of up to $\goods!$ from the search. To
measure its effect we ran the prover at $\agents=4$ for $\goods\in\{6,7\}$ both
with and without the three optimizations, leaving everything else identical;
Table~\ref{tab:bench} reports the wall-clock.

\begin{table}[h]
\centering
\begin{tabular}{rrrr}
\toprule
$\goods$ & \#partitions & with relabelling & without relabelling\\
\midrule
$6$ & $1{,}560$ & $<1$\,s & $8$\,s\\
$7$ & $8{,}400$ & $196$\,s ($\approx 3.3$\,min) & $>30$\,min (did not finish)\\
\bottomrule
\end{tabular}
\caption{Effect of the agent-$1$ relabelling symmetry break ($\agents=4$).
Disabling it leaves the $\goods!$ permutation copies of every profile in the
search, and the solver must rediscover the same refutation for each.}
\label{tab:bench}
\end{table}

The symmetry break helps, and more so as $\goods$ grows. At $\goods=7$ the
unsorted run did not finish within a $30$-minute cap while the sorted run took
about $3$ minutes --- at least a tenfold reduction, and only a lower bound since
the unsorted run did not terminate. We did not run the unsorted $\goods=8$ case.

\section{Results}
\label{sec:results}

For $(\agents,\goods)=(4,8)$ the solver returned \texttt{unsat}
(``\texttt{No counterexample}''), establishing Theorem~\ref{thm:main} for strictly
positive valuations; Section~\ref{sec:nonneg} lifts it to all non-negative
profiles. The
smaller instances $(4,6)$ and $(4,7)$ were solved first and likewise admit a full
\EFX{} allocation for every profile. Table~\ref{tab:results} summarises the runs.

\begin{table}[h]
\centering
\begin{tabular}{lrrl}
\toprule
Instance $(\agents,\goods)$ & \#partitions & Wall-clock & Outcome\\
\midrule
$(4,6)$ & $1{,}560$  & $<1$\,s        & \texttt{unsat} (full \EFX{} exists)\\
$(4,7)$ & $8{,}400$  & $196$\,s       & \texttt{unsat} (full \EFX{} exists)\\
$(4,8)$ & $40{,}824$ & $62$\,h\,$34$\,m & \texttt{unsat} (full \EFX{} exists)\\
\bottomrule
\end{tabular}
\caption{Outcomes of the SMT runs; peak memory for the $(4,8)$ run was about
$415$\,MB.}
\label{tab:results}
\end{table}

\begin{observation}
For this instance size, Theorem~\ref{thm:main} is stronger than the ``almost
full'' guarantee of \cite{BergerCohenFeldmanFiat2022}: it shows that the single
item that their result may leave unallocated is never needed at $(4,8)$ --- a
complete allocation always exists.
\end{observation}

\section{Extension to non-negative valuations}
\label{sec:nonneg}

The solver assumes strictly positive valuations, but the result extends to all
\emph{non-negative} additive profiles, including those with zero entries, and the
argument stays within the additive class.

\begin{theorem}
\label{thm:nonneg}
Fix $\agents,\goods$. If every strictly positive additive profile admits a full
\EFX{} allocation over the ordered partitions, then every
non-negative additive profile admits a full \EFX$_0$ allocation over the same
family.
\end{theorem}

\begin{proof}
The argument perturbs the zero entries to obtain a strictly positive profile,
applies the hypothesis to get a full \EFX{} allocation there (on a strictly
positive profile no good is valued at zero, so this allocation is already
\EFX$_0$), and shows the same allocation stays \EFX$_0$ for the original profile.

Each agent's \EFX$_0$ constraints involve only that agent's own weights, so we
argue one agent $i$ at a time, with weights $w_{i,\cdot}\ge 0$. If all of agent
$i$'s weights are zero, every \EFX$_0$ constraint for $i$ reads $0\ge 0$ and holds;
otherwise $i$ has a positive weight. Let
$W_i=\{\sum_{g\in S}w_{i,g}:S\subseteq[\goods]\}$ be the bundle values $i$ can
attain and set
\[
  d_i=\min\{|x-y|:x,y\in W_i,\ x\neq y\}>0,\qquad
  \varepsilon_i=\frac{d_i}{2\goods}.
\]
(Since $i$ has a positive weight, $W_i$ is a finite set with at least two distinct
elements --- for instance $0$ and $\max_g w_{i,g}>0$ --- so the minimum defining
$d_i$ ranges over a finite, non-empty set of positive gaps and is itself positive.) For each good with $w_{i,g}>0$ the values $0$
and $w_{i,g}$ lie in $W_i$, hence $d_i\le w_{i,g}$ and so
$\varepsilon_i<d_i\le w_{i,g}$: $\varepsilon_i$ is below every positive weight of
$i$.

Let $\hat w$ replace each zero entry $w_{i,g}=0$ by $\varepsilon_i$ and keep
positive entries. Then $\hat w$ is strictly positive, so by hypothesis it admits a
full \EFX{} allocation $X$; we show the same $X$ is
\EFX$_0$ for $w$. Consider a pair $(i,j)$ with $i\neq j$ and let
$g^\star=\arg\min_{g\in X_j}w_{i,g}$; it suffices to check
$v_i(X_i)\ge v_i(X_j\setminus\{g^\star\})$ under $w$. As $\varepsilon_i$ is below
every positive weight, $g^\star$ also minimises $\hat w_{i,\cdot}$ over $X_j$: if
$X_j$ holds a good $i$ values at $0$ that good minimises under both $w$ and
$\hat w$, and otherwise the smallest positive good does. Write $L=v_i(X_i)$,
$R=v_i(X_j\setminus\{g^\star\})$ (values under $w$) and let $a,b$ count the
goods that $i$ values at $0$ in $X_i$ and in $X_j\setminus\{g^\star\}$. The \EFX{} inequality for
$\hat w$ at $g^\star$ is $L+a\varepsilon_i\ge R+b\varepsilon_i$, i.e.
$(a-b)\varepsilon_i\ge R-L$. If \EFX$_0$ under $w$ failed then $L<R$, so $R-L>0$
with $L,R\in W_i$ distinct and hence $R-L\ge d_i$; but $a-b\le a\le\goods$ gives
$(a-b)\varepsilon_i\le\goods\,\varepsilon_i=d_i/2<d_i$, a contradiction. Thus
$L\ge R$, the pair satisfies \EFX$_0$ under $w$, and $X$ is \EFX$_0$ for $w$.
\end{proof}

\begin{corollary}
\label{cor:nonneg}
For each $(\agents,\goods)\in\{(4,6),(4,7),(4,8)\}$, every non-negative additive
profile admits a full \EFX$_0$ allocation in which every agent receives at least
one good.
\end{corollary}

\begin{corollary}[The three guarantees coincide]
\label{cor:equiv}
Fix $\agents,\goods$ and consider full allocations over the ordered partitions of
$[\goods]$ into $\agents$ non-empty bundles. The following are equivalent:
\begin{enumerate}
  \item[\textnormal{(P)}] every strictly positive additive profile admits a full
        \EFX{} allocation;
  \item[\textnormal{(Q)}] every non-negative additive profile admits a full
        \EFX$_0$ allocation;
  \item[\textnormal{(R)}] every non-negative additive profile admits a full
        \EFX$_+$ allocation.
\end{enumerate}
\end{corollary}

\begin{proof}
(P)$\Rightarrow$(Q) is Theorem~\ref{thm:nonneg}. (Q)$\Rightarrow$(R) is immediate:
on any profile the \EFX$_+$ constraints (removing only positively-valued goods) form a
subset of the \EFX$_0$ constraints (removing any good), so every \EFX$_0$ allocation is
\EFX$_+$ and a full \EFX$_0$ allocation is in particular a full \EFX$_+$ one. For
(R)$\Rightarrow$(P), let $w>0$ be strictly positive. It is non-negative, so by (R) it
admits a full \EFX$_+$ allocation $X$ over the ordered partitions; since every entry of
$w$ is positive, no good is valued at zero, so on $w$ the \EFX$_+$, \EFX$_0$ and plain
\EFX{} conditions coincide. Hence $X$ is a full \EFX{} allocation and (P) holds.
\end{proof}

\section{Related work and novelty}
\label{sec:related}

Two lines of work are especially close.

\paragraph{Berger, Cohen, Feldman and Fiat \cite{BergerCohenFeldmanFiat2022}.}
They proved that four additive agents always admit an \EFX{} allocation
\emph{leaving at most one good unallocated}, and more generally that $\agents$
additive agents admit one leaving at most $\agents-2$ goods unallocated. This is
a theorem for \emph{all} $\goods$, but it is an ``almost full'' guarantee:
whether a \emph{complete} four-agent \EFX{} allocation always exists is left
open. Our result removes the leftover item for the specific size $\goods=8$.

\paragraph{Akrami, Mayorov, Mehlhorn, Srinivas and Weidenbach \cite{AkramiEtAl2026}.}
This recent work uses the same reduction idea --- encode the negation of \EFX{}
existence and let \texttt{sat}/\texttt{unsat} decide --- but in a different
parameter regime. Their computations are for \emph{three} agents
($\goods\in\{6,7,8\}$) under \emph{monotone} valuations: they certify existence
for three agents and seven goods, and find by SAT a submodular (hence monotone)
counterexample for three agents and eight goods; the counterexample for general
$n\geq 3,\ m\geq n+5$ is then obtained from this seed by an explicit construction
rather than by solving. Their Z3 SMT-over-LRA encoding did not terminate within
$40$ hours for three agents and seven goods (three agents and six goods took about
$26$ seconds). They explicitly leave the \emph{additive} case open, and their
computational results do not reach four agents.

\begin{table}[h]
\centering
\small
\begin{tabular}{p{3.1cm}cccp{2.3cm}p{2.4cm}}
\toprule
Work & Agents & Goods & Valuations & Guarantee & Method\\
\midrule
Plaut--Roughgarden \cite{PlautRoughgarden2020}
  & $2$ & any & monotone & full \EFX{} exists & proof\\
Chaudhury et al. \cite{ChaudhuryGargMehlhorn2024}
  & $3$ & any & additive & full \EFX{} exists & proof\\
Berger et al. \cite{BergerCohenFeldmanFiat2022}
  & $4$ & any & additive & \emph{almost} full ($\leq 1$ leftover) & proof\\
Akrami et al. \cite{AkramiEtAl2026}
  & $3$ & $7$ & monotone & full \EFX{} exists & SAT\\
Akrami et al. \cite{AkramiEtAl2026}
  & $\geq3$ & $\geq n{+}5$ & monotone & \emph{counterexample} & SAT\,+\,constr.\\
\textbf{This report}
  & $\mathbf{4}$ & $\mathbf{8}$ & \textbf{additive} & \textbf{full \EFX{} exists}
  & \textbf{SMT (LRA)}\\
\bottomrule
\end{tabular}
\caption{The four-agent / eight-good additive cell is, to our knowledge,
established here for the first time.}
\label{tab:related}
\end{table}

\paragraph{Novelty.}
Table~\ref{tab:related} summarises these results. To the best of our knowledge,
full \EFX{} existence for four agents and eight goods under additive valuations
has not been established before, by any method. The two works above are related
but address different questions: Berger et al.\ \cite{BergerCohenFeldmanFiat2022}
give an analytic ``almost full'' guarantee for all sizes but not a complete
allocation, while Akrami et al.\ \cite{AkramiEtAl2026} compute with monotone
valuations for three agents and leave the additive case open. The present result
concerns additive valuations, four agents, and a complete allocation, at the
fixed size $\goods=8$.

Our $(4,8)$ instance sits exactly one good \emph{below} the submodular
counterexample threshold $\goods=\agents+5=9$ of \cite{AkramiEtAl2026}: at
$(4,9)$ a submodular (hence monotone) profile with no \EFX{} allocation provably
exists, whereas the additive case there is open. The natural next case is thus
$(4,9)$: an additive \EFX{} result there would show that additivity, not just
monotonicity, is what makes \EFX{} attainable at this size. We return to it in
Section~\ref{sec:conclusion}.

\section{Discussion and limitations}
\label{sec:discussion}

\paragraph{No machine-checked certificate.}
The $(4,8)$ result rests on Z3 returning \texttt{unsat}, not on an externally
replayable proof object. We attempted to extract one, but enabling Z3's proof
generation did not yield a usable proof: the solver returned \texttt{unknown}
under the proof-producing configuration, so no DRAT- or Lean-checkable certificate
is available. By contrast, Akrami et al.\ \cite{AkramiEtAl2026} verified their SAT
encoding in Lean and their unsatisfiability proofs with \textsc{drat-trim}.
Obtaining a checkable certificate for the $(4,8)$ query would require a different,
proof-producing pipeline (for instance an exact-arithmetic LP/Farkas back-end),
and remains open.

\paragraph{What underpins the trust in the result.}
Two things support Theorem~\ref{thm:main}. (i) The reduction itself is
elementary and is argued in full in Section~\ref{sec:reduction}, with the one
non-obvious step --- restricting to non-empty partitions --- proved sound in
Proposition~\ref{prop:conservative}; the exactness of LRA reasoning
(Section~\ref{sec:lra}) means the \texttt{unsat} verdict is an exact
impossibility over the reals, not a numerical artefact. (ii) The decision is
reproducible: the query is exported as a standalone SMT-LIB~2 script and can be
replayed by any LRA-capable solver, so the result does not depend on our
particular driver code.

\paragraph{Scope of the claim.}
Theorem~\ref{thm:main} covers all
\emph{non-negative additive} valuations at the \emph{fixed} size $\goods=8$. It
says nothing about non-additive valuations, about chores (negative valuations), or
about $\goods=9$ or larger --- each size is a separate, much larger computation we
have not run, and existence at one size does not imply it at the next. The
$\mathcal{P}$-restriction to non-empty bundles
is sound for the existence direction by Proposition~\ref{prop:conservative}, but
means the same code, used to \emph{find} a counterexample, would have to be
extended to allow empty bundles before a \texttt{sat} answer could be trusted as
a valid counterexample.

\paragraph{Scaling.}
The jump from $1{,}560$ to $40{,}824$ partitions ($6\to 8$ goods) took the
runtime from seconds to $62.5$ hours. Each added good both multiplies the
partition count and widens every constraint, so $\goods=9$ is well out of reach
of the present encoding without substantially stronger symmetry breaking or a
fundamentally different decomposition.

\section{Conclusion and future work}
\label{sec:conclusion}

We reduced the existence of full \EFX{} allocations, for a fixed instance size,
to a single satisfiability query over linear real arithmetic, and used it to show
by computer that every additive profile with four agents and eight goods admits a
full \EFX{} allocation (Theorem~\ref{thm:main}). To our knowledge this is the
first verification for the four-agent / eight-good additive case; it strengthens
the ``almost full'' four-agent guarantee of \cite{BergerCohenFeldmanFiat2022} at
this size and complements the recent SAT-based work of \cite{AkramiEtAl2026},
which addressed three agents and left additive valuations open.

Several directions follow naturally:
\begin{itemize}
  \item \textbf{A checkable certificate.} Our attempt to obtain a replayable proof
        failed --- Z3 returned \texttt{unknown} under proof generation
        (Section~\ref{sec:discussion}). A proof-producing pipeline, for instance an
        exact LP/Farkas back-end emitting a certificate per allocation, would
        remove the reliance on trusting the solver.
  \item \textbf{The separating case $\goods=9$.} For $\agents=4$ the smallest
        submodular (hence monotone) \EFX{} counterexample \emph{guaranteed} by
        \cite{AkramiEtAl2026} occurs at $\goods=\agents+5=9$: a submodular profile
        with four agents and \emph{nine} goods provably admits \emph{no} \EFX{}
        allocation.
        Settling the \emph{additive} case at $(4,9)$ is therefore the natural next
        target, and it is informative either way. If the solver
        returns \texttt{unsat}, then additive \EFX{} exists where monotone \EFX{}
        fails --- a \emph{separation} showing that \EFX{} existence depends on
        additivity, in the positive direction of the open question raised by
        \cite{AkramiEtAl2026}. If instead it returns \texttt{sat}, the
        witness is an additive instance with no \EFX{} allocation --- a refutation
        of \EFX{} existence for additive valuations, i.e.\ a resolution of the
        main open problem for additive valuations. The obstacle is scale: $(4,9)$ has
        $4!\,\genfrac{\{}{\}}{0pt}{}{9}{4}=186{,}480$ ordered partitions
        ($\approx 4.6\times$ the $40{,}824$ of $(4,8)$), an extra variable per
        agent, and longer constraints, so it will require parallelising or sharding
        the partition set and symmetry breaking stronger than sorting a single
        agent.
  \item \textbf{Variants.} Adapt the encoding to related notions: \EFX{} for
        chores (negative valuations), or the relaxations (tEFX, and \EFone{}
        together with EEFX) that \cite{AkramiEtAl2026} study for three agents.
\end{itemize}

\printbibliography[heading=bibintoc,title={References}]

\end{document}
